% =====================================================================
%  BAB I PENDAHULUAN
% =====================================================================
\chapter{PENDAHULUAN}
\label{bab:pendahuluan}

\petunjuk{Argumentasi ilmiah yang melatarbelakangi penelitian, arah
penelitian, dan batasannya (Pedoman TA 7.1). \tahap{\WAJIB}{\WAJIB{} (final)}}

\section{Latar Belakang}
\petunjuk{Mulai dari masalah umum di bidang Teknik Informatika lalu mengerucut
ke masalah spesifik. Tunjukkan kesenjangan antara kondisi yang diharapkan dan
kondisi aktual. Wajib didukung data empiris dan rujukan terkini
($\leq$5 tahun).}

\isian{Tulis uraian dalam bentuk paragraf.} Contoh sitasi IEEE: \cite{contoh2024jurnal}.

\section{Rumusan Masalah}
\petunjuk{Pertanyaan penelitian yang jelas, spesifik, terukur, dan dapat
dijawab dengan metode ilmiah. Contoh: Bagaimana kinerja algoritma X
dibandingkan Y pada dataset Z dalam hal akurasi dan waktu eksekusi?}

\begin{enumerate}[leftmargin=1.5em]
  \item \isian{\dots\dots\dots}
  \item \isian{\dots\dots\dots}
\end{enumerate}

\section{Tujuan Penelitian}
\petunjuk{Satu tujuan untuk setiap rumusan masalah; gunakan kata kerja
operasional (menganalisis, merancang, membandingkan, mengevaluasi).
Cantumkan indikator keberhasilan. Uraikan tujuan dan indikator keberhasilan
untuk setiap rumusan masalah.}

\isian{Tulis uraian dalam bentuk paragraf.}

\section{Batasan Penelitian}
\petunjuk{Jenis dan ukuran dataset, platform atau lingkungan eksperimen,
metrik evaluasi, dan hal yang tidak dicakup.}

\isian{Tulis uraian dalam bentuk paragraf.}

\section{Manfaat Penelitian}
\petunjuk{Manfaat teoritis (keilmuan) dan praktis (aplikasi nyata), realistis
dan selaras dengan luaran.}

\isian{Tulis uraian dalam bentuk paragraf.}

\section{Sistematika Penulisan}
\petunjuk{Satu paragraf pendek per bab yang merangkum isi Bab I sampai Bab V.}

\isian{Tulis uraian dalam bentuk paragraf.}
